\section{嵌入子流形}
\label{sec:mrel-embedded-submanifolds}

浸入在每点附近都像把一块平的片放进目标空间，可这些局部的片拼起来，未必是目标中一个像样的子空间：像可能自己穿过自己，也可能从远处绕回来贴近自身。本节先在浸入之上加入整体的拓扑条件，得到嵌入；再在目标的坐标中描述嵌入子流形的局部样子；最后研究由约束方程定义的集合何时是这样的子流形。

\subsection{嵌入}
\label{sec:mrel-embedding}

\paragraph{“8”字形曲线}
考虑$\alpha:(-\pi,\pi)\to\mathbb R^2$，$\alpha(t)=(\sin2t,\sin t)$。速度$(2\cos2t,\cos t)$处处非零：若$\cos t=0$，则$t=\pm\pi/2$，此时$\cos2t=-1$。所以$\alpha$是浸入。它还是单射：若$\sin t=\sin s$且$t\neq s$，则$s=\pi-t$或$s=-\pi-t$，于是$\sin2s=-\sin2t$，要使第一个分量也相等只能$\sin2t=0$，即$t\in\{0,\pm\pi/2\}$。$t=0$时$s=\pm\pi$不在区间内，$t=\pm\pi/2$时$s=t$，都与假设矛盾。特别地，$t=0$是唯一映到原点的参数。

$\alpha$的像是一个“8”字：$t$从$0$增大到$\pi$时画出一个圈，$t\to\pi$时又回到原点附近；$t$从$0$减小到$-\pi$时画出另一个圈。取$t_k\to\pi$，则$\alpha(t_k)\to(0,0)=\alpha(0)$，但$t_k$不趋于$0$。从像的角度看，点列收敛到原点；从参数的角度看，对应的参数却跑到了区间的另一端。所以$\alpha$作为从$(-\pi,\pi)$到其像的双射，逆映射不连续。

\paragraph{定义}
“8”字形曲线的问题不在局部方向，而在像所继承的拓扑与源流形的拓扑不一致。像$F(M)$作为$N$的子集，有从$N$继承的子空间拓扑：$F(M)$中的开集是$N$中开集与$F(M)$的交。要求$F$在这一拓扑下是同胚，就排除了上述情形。

\begin{definition}[嵌入]
\label{def:mrel-embedding}
光滑映射$F:M\to N$称为\term{嵌入}（embedding），若$F$是浸入，并且$F:M\to F(M)$是同胚，其中$F(M)$带有从$N$继承的子空间拓扑。
\end{definition}

于是有三个层层加强的条件：浸入只要求微分单射；单射浸入还要求映射本身单射；嵌入进一步要求到像的同胚。$E:\mathbb R\to\mathbb R^2$，$t\mapsto(\cos t,\sin t)$是浸入但不是单射，像自己反复重叠；$\alpha$是单射浸入但不是嵌入；把$E$限制到$(0,2\pi)$上，得到单射浸入，它的逆就是角度图$\theta_2$，连续，所以是嵌入。

\paragraph{紧流形的单射浸入}
源流形紧时，后两个条件之间的差距消失。

\begin{theorem}[紧流形上的单射浸入]
\label{thm:mrel-compact-injective-immersion}
若$M$紧，$F:M\to N$是单射浸入，则$F$是嵌入。
\end{theorem}

\begin{proof}
只需证明$F:M\to F(M)$的逆连续，即$F$把$M$的闭集映为$F(M)$的闭集。设$K\subseteq M$是闭集，紧空间的闭子集是紧集，所以$K$紧；连续映射把紧集映为紧集，$F(K)$紧；$N$是Hausdorff空间，其紧子集是闭集，所以$F(K)$在$N$中闭，从而在$F(M)$中闭。
\end{proof}

$\alpha$的源区间$(-\pi,\pi)$不紧，参数可以逃向端点而像却回到原点，定理正是排除了这种逃逸。

\subsection{子流形的局部表示}
\label{sec:mrel-submanifold-local}

\paragraph{切片坐标}
嵌入的像在目标中应当局部像一块平的片。用目标的坐标表达：在合适的图中，像恰好是一部分坐标为零的那片点。

\begin{definition}[嵌入子流形]
\label{def:mrel-embedded-submanifold}
设$M$是$n$维光滑流形，$S\subseteq M$，$0\leq k\leq n$。若对每个$p\in S$，都有$M$的图$(U,\varphi)$，$p\in U$，使
\begin{equation}
\varphi(U\cap S)=\{\bm x\in\varphi(U):x_{k+1}=\cdots=x_n=0\},
\label{eq:mrel-slice-chart}
\end{equation}
则称$S$为$M$的$k$维\term{嵌入子流形}（embedded submanifold）。这样的图称为$S$的\term{切片坐标}（slice chart）。$n-k$称为$S$的\term{余维数}（codimension）。
\end{definition}

切片坐标中的前$k$个坐标在$S$上自由变化，是$S$的内部坐标；后$n-k$个坐标描述离开$S$的方向，在$S$上为零。这一描述与嵌入是同一件事的两面。

\begin{theorem}[嵌入与嵌入子流形]
\label{thm:mrel-embedding-submanifold}
\begin{enumerate}
\item 设$S$是$M$的$k$维嵌入子流形，带子空间拓扑。把切片坐标限制到$U\cap S$并只保留前$k$个坐标，得到$S$的光滑图册，$S$成为$k$维光滑流形，包含映射$\iota:S\to M$是嵌入。
\item 若$F:P\to M$是嵌入，则$F(P)$是$M$的嵌入子流形，维数等于$\dim P$，并且$F:P\to F(P)$是微分同胚。
\end{enumerate}
\end{theorem}

\begin{proof}
第一条。限制后的映射是$U\cap S$到$\mathbb R^k$开集的同胚，因为$\varphi$是同胚，而式\eqref{eq:mrel-slice-chart}说明$U\cap S$恰好对应坐标平面$\mathbb R^k\times\{\bm 0\}$与$\varphi(U)$的交。两张这样的图之间的坐标变换，是$M$的坐标变换限制到$\mathbb R^k\times\{\bm 0\}$上再取前$k$个分量，仍光滑。$S$作为$M$的子空间是Hausdorff且第二可数的。在这些图中，$\iota$的坐标表示是$\bm x'\mapsto(\bm x',\bm 0)$，是浸入；$S$本来就带子空间拓扑，$\iota$到像的同胚是自然的。

第二条。设$p\in P$，$k=\dim P$。$F$是浸入，由常秩定理，有以$p$为中心的图$(U,\varphi)$与以$F(p)$为中心的图$(V,\psi)$，使$\psi\circ F\circ\varphi^{-1}(\bm x')=(\bm x',\bm 0)$。这只说明$F(U)$是一片，还须排除$F(P)$的其他部分从远处进入$V$。这就要用到同胚条件：$F(U)$在$F(P)$中开，所以存在$M$中的开集$W$，使$F(U)=W\cap F(P)$。

令$V'=V\cap W\cap\psi^{-1}\bigl(\varphi(U)\times\mathbb R^{n-k}\bigr)$。若$q\in V'\cap F(P)$，则$q\in W\cap F(P)=F(U)$，所以$\psi(q)=(\bm x',\bm 0)$，$\bm x'\in\varphi(U)$。反过来，若$q\in V'$且$\psi(q)=(\bm x',\bm 0)$，则$\bm x'\in\varphi(U)$，$q=\psi^{-1}(\bm x',\bm 0)=F(\varphi^{-1}(\bm x'))\in F(P)$。所以$\psi$限制到$V'$上是$F(P)$的切片坐标。最后，在$P$的图$\varphi$与$F(P)$的限制图下，$F$的坐标表示是恒等映射，故$F$是到像的微分同胚。
\end{proof}

\paragraph{子流形的切空间}
包含映射$\iota$是浸入，$\mathrm d\iota_p$把$T_pS$单射地送入$T_pM$。今后把$T_pS$与它的像$\mathrm d\iota_p(T_pS)$等同，看作$T_pM$的$k$维子空间。在切片坐标中，$\iota$的坐标表示是$\bm x'\mapsto(\bm x',\bm 0)$，所以
\[
T_pS=\operatorname{span}\bigl\{\partial_1|_p,\ldots,\partial_k|_p\bigr\}\subseteq T_pM .
\]
沿后$n-k$条坐标线运动会离开$S$，相应的方向不在$T_pS$中。当$M=\mathbb R^n$时，$T_pM$就是$\mathbb R^n$，$T_pS$成为$\mathbb R^n$的一个线性子空间，这就是嵌入在欧氏空间中的曲线与曲面的切向量的通常含义。

\paragraph{例子：圆周、球面、图像与乘积}
圆周$S^1\subseteq\mathbb R^2$。在去掉负半轴$\{(y_1,0):y_1\leq0\}$的平面上用极坐标$(\theta,r)$，令$\varphi(\bm y)=(\theta,r-1)$。在这张图中，$S^1$恰好对应第二个坐标为零，所以它是切片坐标；另一张图用去掉正半轴的极坐标。于是$S^1$是$\mathbb R^2$的一维嵌入子流形，余维数为$1$。三种坐标在这里分得很清楚：$\theta$是圆周的内部坐标，$(y_1,y_2)$是环境坐标，$r-1$是只在环境中有意义的离开圆周的坐标。球面$S^2\subseteq\mathbb R^3$用球坐标$(\theta,\phi,r-1)$类似处理。

光滑函数$f:W\to\mathbb R^{n-k}$（$W\subseteq\mathbb R^k$开）的图像$\Gamma_f=\{(\bm x,f(\bm x)):\bm x\in W\}$是$W\times\mathbb R^{n-k}$的嵌入子流形：$\varphi(\bm x,\bm y)=(\bm x,\bm y-f(\bm x))$是微分同胚，逆为$(\bm x,\bm z)\mapsto(\bm x,\bm z+f(\bm x))$，它把$\Gamma_f$送到$\{\bm z=\bm0\}$。这个例子说明，局部上“是图像”与“有切片坐标”是一回事，下一小节的隐函数表示正是把子流形局部写成图像。

乘积也保持子流形关系。若$S_1\subseteq M_1$与$S_2\subseteq M_2$分别是嵌入子流形，两张切片坐标的乘积是$S_1\times S_2$在$M_1\times M_2$中的切片坐标，只需把坐标重新排序，使取零的坐标放在最后。例如环面$S^1\times S^1$是$\mathbb R^2\times\mathbb R^2=\mathbb R^4$的二维嵌入子流形，余维数为$2$。

\paragraph{不是子流形的子集}
子流形的定义对每个点都要求切片坐标，任意子集未必满足。十字形$X=\{(y_1,y_2):y_1y_2=0\}$在原点以外是两条坐标轴，局部是直线；但在原点处，它不是一维子流形。若是，原点在$X$中有一个连通邻域同胚于开区间，去掉原点后恰有两个连通分支；可原点在$X$中的任何足够小的连通邻域，去掉原点后都有四个分支，即四条半轴的片段。“8”字形曲线的像在原点处也是同样的十字形交叉，所以$\alpha$的像不是嵌入子流形，与$\alpha$不是嵌入相符。

还有更隐蔽的情形。尖点曲线$\{y_2^2=y_1^3\}$作为拓扑空间同胚于实轴，没有十字交叉，但可以证明它在原点处没有切片坐标，因而不是$\mathbb R^2$的光滑嵌入子流形，这里不给出证明。上一节的参数化$\beta(t)=(t^2,t^3)$恰好在原点速度为零，与此相呼应。可见判断是否为子流形，既要看拓扑，也要看光滑结构。

\subsection{正则水平集}
\label{sec:mrel-regular-level-sets}

很多子流形由约束方程给出：球面是$\lVert\bm y\rVert_2^2=1$的解集，正交矩阵是$\bm A^\top\bm A=\bm I$的解集。问题是：约束方程的解集何时是子流形，维数是多少，切空间是什么？

\paragraph{正则值}
设$\Phi:M\to N$光滑，$c\in N$。水平集$\Phi^{-1}(c)$由所有满足约束$\Phi(p)=c$的点组成。若每个约束都提供一个独立的方向限制，就应当把$m$个自由度减去$n$个。“独立”的准确含义是微分满秩。

\begin{definition}[正则值与正则水平集]
\label{def:mrel-regular-value}
设$\Phi:M\to N$光滑。若$\mathrm d\Phi_p$是满射，称$p$为$\Phi$的正则点，否则称为临界点。若$\Phi^{-1}(c)$中的每一点都是正则点，称$c$为$\Phi$的\term{正则值}（regular value），$\Phi^{-1}(c)$称为\term{正则水平集}（regular level set）。
\end{definition}

按定义，不在像中的点也是正则值，这时水平集为空，下面的定理对空集自然成立。

\begin{theorem}[正则水平集定理]
\label{thm:mrel-regular-level-set}
设$\Phi:M\to N$光滑，$c$是$\Phi$的正则值。则$S=\Phi^{-1}(c)$是$M$的$m-n$维嵌入子流形，并且对每个$p\in S$，
\[
T_pS=\ker\mathrm d\Phi_p .
\]
\end{theorem}

\begin{proof}
正则点组成的集合$M_0$是开集，因为满秩在附近保持。$S\subseteq M_0$，$\Phi$在$M_0$上是淹没，由常秩定理，每个$p\in S$附近有图$\varphi$与以$c$为中心的图$\psi$，使$\psi\circ\Phi\circ\varphi^{-1}(x_1,\ldots,x_m)=(x_1,\ldots,x_n)$。于是$S$在这张图中恰好是$x_1=\cdots=x_n=0$的点。把坐标重新排序，让这$n$个坐标排在最后，就得到$m-n$维的切片坐标。

$\Phi$在$S$上取常值$c$，所以$\Phi\circ\iota$是常值映射，微分为零，即$\mathrm d\Phi_p\circ\mathrm d\iota_p=0$，$T_pS\subseteq\ker\mathrm d\Phi_p$。$\mathrm d\Phi_p$满射，由秩--零度关系，核的维数是$m-n$，与$T_pS$相同，所以二者相等。
\end{proof}

定理与隐函数定理是同一件事的两种说法。在欧氏情形，$\Phi:\mathbb R^m\to\mathbb R^n$在$p$处的Jacobian矩阵行满秩，可以选出$n$列组成可逆方阵；隐函数定理把对应的$n$个变量解成其余$m-n$个变量的光滑函数，水平集在$p$附近就是这个函数的图像，再由上一小节的图像例子得到切片坐标。满秩条件排除了局部奇点：在不满秩的点，约束之间可能互相重复，也可能在该点退化，解集在那里可以出现交叉或孤立点。

\paragraph{球面}
令$\Phi:\mathbb R^3\to\mathbb R$，$\Phi(\bm y)=\lVert\bm y\rVert_2^2$。$\mathrm d\Phi_{\bm y}(\bm v)=2\bm y^\top\bm v$，当$\bm y\neq\bm 0$时它是到$\mathbb R$的满射。$\Phi^{-1}(1)=S^2$不含原点，所以$1$是正则值，$S^2$是二维嵌入子流形，并且
\[
T_{\bm p}S^2=\{\bm v\in\mathbb R^3:\bm p^\top\bm v=0\},
\]
即与$\bm p$正交的平面。这与第~\ref{chap:manifold-basics}章用球极投影算出的切向量$(0.2,0,0)$与$\bm p=(0,\tfrac35,\tfrac45)$正交一致。同样，椭球面是$\sum_iy_i^2/a_i^2=1$的正则水平集；线性映射$\bm y\mapsto(y_1/a_1,y_2/a_2,y_3/a_3)$把它一一映到球面，在两侧的子流形结构下是微分同胚，这就是本章开头说椭球面与球面“本质上相同”的含义。

\paragraph{正则性只是充分条件}
不满足正则性时，水平集可能不是子流形。$\Phi(y_1,y_2)=y_1^2-y_2^2$在原点处微分为零，$\Phi^{-1}(0)$是两条直线$y_2=\pm y_1$的并，在原点处交叉，与上一小节的十字形一样不是子流形。$\Phi(\bm y)=\lVert\bm y\rVert_2^2$的零水平集只有原点，是零维子流形，却不是定理预测的二维；那里原点是临界点，定理本来就不适用。

水平集也可能在正则性失败时仍是子流形。$\Phi(\bm y)=(\lVert\bm y\rVert_2^2-1)^2$的零水平集仍是单位圆，但在圆上$\mathrm d\Phi=4(\lVert\bm y\rVert_2^2-1)\bm y^\top=\bm 0$，没有一个点是正则点。这说明正则值是水平集为子流形的充分条件，不是必要条件；同一个集合换一个约束函数，正则性可能成立，也可能不成立。

\paragraph{正交矩阵与重复方程}
全体$n$阶正交矩阵组成的集合$\mathrm O(n)=\{\bm A\in\mathbb R^{n\times n}:\bm A^\top\bm A=\bm I\}$称为\term{正交群}（orthogonal group）。它由$n^2$个方程定义，但这些方程有重复：$\bm A^\top\bm A$总是对称矩阵，第$(i,j)$个方程与第$(j,i)$个方程相同。若把约束映射看成取值于$\mathbb R^{n\times n}$，它的微分只能取到对称矩阵，永远不是满射，单位矩阵不是正则值，定理用不上。

正确的做法是让约束映射取值于对称矩阵组成的空间$\operatorname{Sym}(n)$，它是$n(n+1)/2$维线性空间。令$\Phi:\mathbb R^{n\times n}\to\operatorname{Sym}(n)$，$\Phi(\bm A)=\bm A^\top\bm A$。对方向$\bm H$求导，
\[
\mathrm d\Phi_{\bm A}(\bm H)=\bm H^\top\bm A+\bm A^\top\bm H .
\]
在$\bm A\in\mathrm O(n)$处，任给对称矩阵$\bm C$，取$\bm H=\tfrac12\bm A\bm C$，利用$\bm A^\top\bm A=\bm I$与$\bm C^\top=\bm C$，得
\[
\mathrm d\Phi_{\bm A}\bigl(\tfrac12\bm A\bm C\bigr)=\tfrac12\bm C^\top\bm A^\top\bm A+\tfrac12\bm A^\top\bm A\bm C=\bm C .
\]
所以$\mathrm d\Phi_{\bm A}$是满射，$\bm I$是正则值。

由定理~\ref{thm:mrel-regular-level-set}，$\mathrm O(n)$是$\mathbb R^{n\times n}$的嵌入子流形，维数为$n^2-n(n+1)/2=n(n-1)/2$。在单位矩阵处，$T_{\bm I}\mathrm O(n)=\ker\mathrm d\Phi_{\bm I}=\{\bm H:\bm H^\top+\bm H=\bm 0\}$，即全体反对称矩阵，维数恰为$n(n-1)/2$，与上面一致。$n=2$时维数为$1$，对应平面旋转的一个转角参数；$n=3$时维数为$3$，对应三维旋转的三个自由度。$\mathrm O(n)$既是流形又是群，这种结构在第~\ref{chap:lie-groups}章研究。

\paragraph{带边界的情形}
本章的结论都针对无边界流形。约束若是不等式，例如闭球$\{\bm y:\lVert\bm y\rVert_2^2\leq1\}$，所得集合是带边界流形，边界恰是正则水平集$\{\lVert\bm y\rVert_2^2=1\}$。一般地，若$c$是光滑实值函数$\Phi:M\to\mathbb R$的正则值，下水平集$\{\Phi\leq c\}$是与$M$同维的带边界流形，边界为$\Phi^{-1}(c)$；证明要在边界点处单独构造半空间坐标，这里不给出。约束集合本身带边界，或映射的定义域带边界时，上述定理都需要另核条件。
